% Paper 5, §2 — What "measuring an LLM" means
% Pass-C section. Fixes the eval-methodology vocabulary the rest of the
% paper inherits: the score random variable $\hat{S}(\theta,\mathcal{B},\pi,\tau)$,
% the benchmark instance $\mathcal{B}=(\rho, A, V)$, the three orthogonal
% measurement axes, validity/reliability/construct/contamination/ecological
% validity, the IRT 2PL form, and the HELM scenario-by-metric matrix.
% Anchors: kb/notes/evaluation/eval-methodology.md, kb/excerpts/liang2022-helm.md.
% Forward-refs: §3 (\cref{sec:knowledge-benchmarks}) inherits $(\rho, A, V)$;
% §4 (\cref{sec:contamination}) inherits $\hat{S}(\theta,\mathcal{B},\pi,\tau)$;
% §7-§13 (alignment-threat sub-layers) inherit the capability-vs-alignment axis.

\section{What ``measuring an LLM'' means}
\label{sec:measuring-framing}
\label{sec:eval-methodology}

A benchmark score is a noisy random variable, not a property of the
model. Treating it as a property is the methodological mistake the
2024--2026 field has worked hardest to undo, and naming what is
actually being measured --- against what construct, with what
verifier, under what prompting protocol, on what subset of the
deployment distribution --- is the precondition for everything that
follows in this paper. \textbf{The thesis of this section is that the
2024--2026 measurement landscape is organized by three orthogonal
axes (capability vs alignment; knowledge vs reasoning vs agency;
static benchmark vs adversarial probe vs deployment audit), tied
together by a small psychometric vocabulary (validity, reliability,
construct, contamination, ecological validity), and that the closest
the field has to a theory of measurement is HELM's
scenario-by-metric matrix \citep{liang2022-helm} together with the
2PL form from item response theory.} The vocabulary fixed here is
inherited by \cref{sec:knowledge-benchmarks,sec:contamination,sec:reasoning-benchmarks,sec:agentic-benchmarks,sec:adversarial-robustness,sec:dangerous-capability,sec:sycophancy,sec:scheming,sec:alignment-faking,sec:watermarking,sec:scalable-oversight}
without re-derivation; what changes section to section is what
$(\rho, A, V)$ contains and what construct $\hat{S}$ is supposed to
measure.

\subsection{The score is a noisy random variable}
\label{sec:measuring-score-rv}

Following \citet{liang2022-helm} and the open-source eval-harness
discipline --- in particular EleutherAI's
\texttt{lm-evaluation-harness} \citep{lm-eval-harness}, the
de-facto reference implementation for the academic-benchmark
column of \cref{sec:knowledge-benchmarks}\footnote{KB excerpt:
\texttt{kb/excerpts/lm-eval-harness.md\#sec-framing}.} ---
write a reported benchmark score as a sample from a distribution%
\footnote{KB note: \texttt{kb/notes/evaluation/eval-methodology.md\#1-formal-framing-the-score-is-a-noisy-random-variable}.}
\begin{equation}
\hat{S}(\theta, \mathcal{B}, \pi, \tau) \in [0, 1]
\label{eq:score-rv}
\end{equation}
where $\theta$ is the model under evaluation, $\mathcal{B}$ is the
benchmark instance (item corpus together with verifier $V$), $\pi$ is
the prompt formatting protocol (system prompt, few-shot examples,
answer-extraction template, scoring request type), and $\tau$ is the
decoding configuration (temperature, top-$p$, max tokens, sampling vs
greedy). The reported number $\hat{S}$ is a single draw from this
distribution. The methodological project is twofold: bound the
variance of $\hat{S}$ across $(\pi, \tau)$, and detect when
$\mathcal{B}$ has been compromised by training-set leakage so the
sample mean of $\hat{S}$ is biased upward by an amount that depends
on memorization strength rather than capability.

The benchmark instance itself unpacks differently per layer of
\cref{sec:introduction}. For the knowledge and reasoning layers
(\cref{sec:knowledge-benchmarks,sec:reasoning-benchmarks}) the
instance is a triple
\begin{equation}
\mathcal{B} = (\rho, A, V)
\label{eq:bench-static}
\end{equation}
where $\rho$ is the problem statement, $A$ is the answer-format
constraint (one of $\{$A,B,C,D$\}$, an integer in $[0, 999]$, a
parsed symbolic expression, a short string), and
$V(\hat{a}, a^\star) \in \{0, 1\}$ is the verifier on the model's
chosen answer $\hat{a}$ against the gold answer $a^\star$. There is
no environment, no trajectory, no end-state. For the agentic layer
\cref{sec:agentic-benchmarks} expands the instance to include initial
state, action set, transition, and trajectory-grading evaluator; for
the adversarial-robustness layer \cref{sec:adversarial-robustness}
adds an attacker policy and an attack budget. The static-instance
form \eqref{eq:bench-static} is the simplest case and the one to
which subsequent sections add structure.

\subsection{Three failure modes that compound}
\label{sec:measuring-failure-modes}

Three failure modes break the cross-paper, cross-lab comparability of
$\hat{S}$, and they do so additively in expectation%
\footnote{KB note: \texttt{kb/notes/evaluation/eval-methodology.md\#1-formal-framing-the-score-is-a-noisy-random-variable}.}.

\paragraph{Prompt-format sensitivity.} A single benchmark item can be
presented many ways: loglikelihood scoring on
$\{\log P(\text{A}), \log P(\text{B}), \log P(\text{C}), \log P(\text{D})\}$
with argmax extraction; generation scoring with regex parsing of
``the answer is X.''; system-prompt presence; zero- vs five- vs
twenty-five-shot in context. Each moves $\hat{S}$ by approximately
$1$--$4$ percentage points on MMLU; combined effects can exceed $10$.
\citet{wang2024-mmlu-pro} test $24$ prompt styles per model and report
prompt-induced score variance falling from approximately $4$--$5\%$ on
MMLU to approximately $2\%$ on MMLU-Pro
\citep[\S{}abstract]{wang2024-mmlu-pro}, the latter achieved by
$10$-way (rather than $4$-way) MCQ that dilutes the answer-letter
prior. The implication: $\pi$ in \eqref{eq:score-rv} is a benchmark
axis, not a hygiene afterthought.

\paragraph{Contamination.} Items in $\mathcal{B}$ may be in $\theta$'s
training set verbatim, near-verbatim, or in distribution
\citep{contamination-survey-2025}. This biases $\hat{S}$ upward by an
amount that depends on memorization strength. \cref{sec:contamination}
takes contamination as its own layer; for the framing here it
suffices that contamination is a property of the pair
$(\theta, \mathcal{B})$, not of either alone, and the bias can be
positive even when no benchmark item appears verbatim in pretraining
(distributional and indirect contamination remain).

\paragraph{Verifier mis-specification.} $V$ in \eqref{eq:bench-static}
may accept wrong answers or reject correct ones.
\citet{mmlu-redux-2024} find $6.49\%$ of MMLU questions contain
errors --- malformed questions, wrong gold answers, ambiguous wording
--- that propagate as a hard upper bound on the perfectly-calibrated
true score on gold-labeled MMLU; the Virology subset contains errors
in $57\%$ of audited questions, which means the per-subject score
there is essentially verifier noise. Verifier error and contamination
compound: a model that has memorized the \emph{wrong-but-popular}
gold answer outscores a model with better true reasoning, because the
wrong popular answers preferentially appear in the same training
corpora the models were trained on.

\subsection{Validity, reliability, construct, ecological validity}
\label{sec:measuring-validity}

The psychometric vocabulary the field borrows --- and most often
borrows incompletely --- decomposes the methodological problem into
four named questions, each of which has a partial 2024--2026 answer.

\begin{description}
  \item[Reliability.] Across independent runs of the same evaluation
    on the same model, do we get the same number? Reliability is what
    bounded $\pi$-and-$\tau$ variance buys: it is necessary for any
    cross-model comparison. \citet{wang2024-mmlu-pro}'s
    $24$-prompt-style protocol is a reliability instrument; the
    $4\text{--}5\%\to 2\%$ drop is a reliability improvement, not a
    capability one \citep[\S{}abstract]{wang2024-mmlu-pro}.

  \item[Validity.] Does the score reflect the construct it claims to
    measure? A score on a contaminated benchmark is reliable but not
    valid for the construct ``capability on unseen items.'' A score
    on a verifier-broken benchmark (e.g., MMLU Virology) is reliable
    but not valid for the construct ``knows virology.'' Validity
    concerns dominate the contradictions chapter
    (\cref{sec:contradictions}).

  \item[Construct.] What latent property is the benchmark trying to
    expose? The construct of MMLU is broad-but-shallow factual
    recall; of FrontierMath, research-grade mathematical problem
    solving; of HarmBench, refusal-rate under attack; of Apollo's
    scheming suite, in-context capability for strategic deception.
    A frequent failure mode is construct drift: a benchmark whose
    items were chosen to expose construct $C_1$ ends up being
    reported as evidence about construct $C_2$, where the item
    distribution does not in fact license that inference.

  \item[Ecological validity.] Does the benchmark setup look anything
    like the deployment distribution? Static MCQ on academic exams
    has low ecological validity for deployment as a general
    assistant; agentic benchmarks
    (\cref{sec:agentic-benchmarks}) trade ecological validity
    against measurability, and that trade is the layer's defining
    failure mode.
\end{description}

The four are not independent. A reliability fix (MMLU-Pro's
$24$-prompt protocol) does not buy validity if the construct is wrong
(saturated retrieval). A validity fix (FrontierMath's unpublished
items) does not buy ecological validity if the deployment distribution
is web-search-augmented dialogue, not closed-book exam. Naming the
four lets us say cleanly which problems each layer's instruments do
and do not address.

\subsection{Three orthogonal measurement axes}
\label{sec:measuring-axes}

The body sections compose against three orthogonal axes. Listing them
here lets each subsequent section locate itself.

\paragraph{Axis 1 --- capability vs alignment.} What the model
\emph{can} do versus what it \emph{will} do under deployment
conditions. Knowledge benchmarks (\cref{sec:knowledge-benchmarks}),
reasoning benchmarks (\cref{sec:reasoning-benchmarks}), agentic
benchmarks (\cref{sec:agentic-benchmarks}), and dangerous-capability
evaluations (\cref{sec:dangerous-capability}) live on the capability
side: they ask whether the model has a capacity at all. Sycophancy
(\cref{sec:sycophancy}), in-context scheming
(\cref{sec:scheming}), and alignment-faking
(\cref{sec:alignment-faking}) live on the alignment side: they ask
whether the model behaves aligned because it \emph{is} aligned, or
because it has the capability and incentive to look aligned during
evaluation. Adversarial robustness (\cref{sec:adversarial-robustness})
sits at the boundary --- a jailbroken model is a model whose
deployment behavior diverges from its trained behavior under search
pressure. The capability-vs-alignment split is the
single most load-bearing axis of the paper, and it pre-stages the
\cref{sec:sycophancy,sec:scheming,sec:alignment-faking} progression
where the failure mode requires successively more strategic
cognition on the model's part.

\paragraph{Axis 2 --- knowledge vs reasoning vs agency.} Recall
versus multi-step inference versus end-to-end task completion in an
environment. The split is empirically diagnostic: in MMLU-Pro,
chain-of-thought scoring meaningfully outperforms direct-answer
scoring on the new harder items, where on the original MMLU it did
not \citep[\S{}sec-cot]{wang2024-mmlu-pro}. The CoT-vs-direct gap
is the operational signature of reasoning-dominant items; a benchmark
where it is zero is dominated by retrieval. Agentic benchmarks
(\cref{sec:agentic-benchmarks}) extend the axis again by adding an
environment whose state and tool affordances the model must navigate.

\paragraph{Axis 3 --- static benchmark vs adversarial probe vs
deployment audit.} Fixed-set scoring versus red-team-style search
versus observational measurement during real use. Static benchmarks
(\cref{sec:knowledge-benchmarks,sec:reasoning-benchmarks}) hold
$\mathcal{B}$ and $\pi$ fixed. Adversarial probes
(\cref{sec:adversarial-robustness}) introduce an attacker that
searches over inputs $\rho$ to maximize $\Pr[\text{policy violation}]$;
the score is attack success rate on a search-augmented set, not
pass-rate on a static one. Deployment audits
(\cref{sec:scheming,sec:alignment-faking,sec:watermarking}) are
observational; $\mathcal{B}$ is not fixed by the evaluator and $\pi$
is whatever the user typed. The 2026 field has plenty of axis-1 and
axis-2 instruments and very few axis-3 instruments; the gap is one
of the contradictions catalogued in \cref{sec:contradictions}.

\begin{intuition}
The three axes are orthogonal in the sense that any
benchmark-as-measurement-instrument occupies a triple
$(\text{capability/alignment}, \text{knowledge/reasoning/agency},
\text{static/adversarial/audit})$ and changing one coordinate
generally requires a different instrument. MMLU is
$(\text{capability}, \text{knowledge}, \text{static})$; HarmBench is
$(\text{alignment-adjacent}, \text{agency-adjacent},
\text{adversarial})$; the (yet-unbuilt) deployment audit for
in-the-wild scheming would be
$(\text{alignment}, \text{agency}, \text{audit})$. Returning to
canonical form: the score random variable
$\hat{S}(\theta, \mathcal{B}, \pi, \tau)$ takes its support, its
verifier, and its admissible $(\pi, \tau)$ ranges from a different
distribution at each cell of the $2\times 3\times 3 = 18$ axis grid,
and there is no single benchmark that spans the grid.
\end{intuition}

\subsection{HELM's scenario-by-metric matrix}
\label{sec:measuring-helm}

HELM \citep{liang2022-helm} is the closest the field has to a
theory of measurement at the scenario level. The structural argument
is that single-axis benchmarks (accuracy alone) systematically miss
the failure modes that determine deployment behavior, and that a
benchmark instance should ship a metric \emph{bundle} rather than a
single number%
\footnote{KB excerpt: \texttt{kb/excerpts/liang2022-helm.md\#sec-metrics}.}.
HELM operationalizes this as a scenario-by-metric matrix: $7$ metrics
$\times$ $16$ core scenarios, with measurements taken
``for each of $16$ core scenarios when possible ($87.5\%$ of the
time)'' \citep[\S{}abstract]{liang2022-helm}.

The seven metric axes are
\citep[\S{}sec-metrics]{liang2022-helm}:
\begin{equation}
\mathcal{M} = \{
  \text{Acc}, \text{Cal}, \text{Rob}, \text{Fair}, \text{Bias},
  \text{Tox}, \text{Eff}
\}
\label{eq:helm-metrics}
\end{equation}
where Acc is task-specific accuracy (exact match, F1, BLEU); Cal is
calibration (Brier or expected calibration error between predicted
probability and empirical accuracy); Rob is robustness (accuracy
delta under perturbed prompts); Fair is the accuracy gap across
demographic-subgroup-conditioned prompts; Bias is demographic
representation in open-ended generation; Tox is toxicity rate of
generations under a fixed classifier; Eff is wall-clock and energy
per inference call. A scenario $s$ is a triple
$(\rho_s, a^\star_s, V_s)$, and the HELM evaluation of model
$\theta$ on scenario $s$ produces a vector
\begin{equation}
\hat{\mathbf{S}}(\theta, s) =
  \big(\hat{S}_{\text{Acc}}, \hat{S}_{\text{Cal}}, \hat{S}_{\text{Rob}},
       \hat{S}_{\text{Fair}}, \hat{S}_{\text{Bias}},
       \hat{S}_{\text{Tox}}, \hat{S}_{\text{Eff}}\big)(\theta, s)
\label{eq:helm-vector}
\end{equation}
not a scalar. The scenario-by-metric matrix is the outer product of
\eqref{eq:helm-metrics} with the $16$-scenario set, populated where
the metric is well-defined for the scenario%
\footnote{KB excerpt: \texttt{kb/excerpts/liang2022-helm.md\#sec-metrics}.}.

The headline coverage finding is the methodological contribution:
``Before HELM, models on average were evaluated on just $17.9\%$ of
the core HELM scenarios''; HELM evaluates $30$ prominent language
models on all $42$ scenarios under a single harness
\citep[\S{}sec-coverage]{liang2022-helm}. The reproducibility move
is to release all raw model prompts and completions, not just
aggregate scores \citep[\S{}sec-output-discipline]{liang2022-helm}.
The metric bundle propagated only partially to subsequent frameworks
--- 2024 leaderboards remain single-metric-dominant --- but the
prompt-and-completion artifact discipline is now standard.

\begin{contradiction}
The HELM proposal is that $\hat{S}$ is a vector
\eqref{eq:helm-vector}, not a scalar. The field's actual reporting
practice in $2024$--$2026$ remains overwhelmingly scalar: model cards
publish a single MMLU number, a single GPQA number, a single
HarmBench ASR. Per-axis breakdowns appear in tech reports
intermittently and rarely with the seven-metric coverage HELM
prescribes. The methodological proposal is uncontroversial; the
adoption gap is real, and several of the contradictions catalogued
in \cref{sec:contradictions} are downstream of this gap (a single
``ASR'' number across HarmBench, JailbreakBench, AdvBench is
exactly the kind of scalarization HELM warns against).
\end{contradiction}

\subsection{Item response theory: a calibration tool}
\label{sec:measuring-irt}

Where HELM addresses the metric-bundle question, item response theory
(IRT) addresses the within-metric calibration question: under a
fixed metric (say accuracy on multiple-choice items), what is the
probability that a model of latent ability $\theta$ answers a
particular item correctly, as a function of that item's difficulty
$b$ and discrimination $a$? The two-parameter logistic (2PL) form
\deepencite{lord1980-irt-applications §1, citation-pending} is
\begin{equation}
P(\text{correct} \mid \theta, b, a) = \sigma\!\big(a(\theta - b)\big)
  = \frac{1}{1 + \exp(-a(\theta - b))}
\label{eq:irt-2pl}
\end{equation}
where $\sigma$ is the logistic, $\theta$ is the model's latent
ability on the construct the item set is supposed to measure, $b$ is
the item difficulty (the ability level at which the item is
answered correctly with probability $1/2$), and $a$ is the item's
discrimination (how steeply the response curve rises across
$\theta$). Items with high $a$ separate models cleanly; items with
low $a$ are close to noise.

The use of \eqref{eq:irt-2pl} in LLM evaluation is calibration, not
benchmark construction. Given a benchmark $\mathcal{B}$ with item
parameters $\{(b_i, a_i)\}$ estimated from a panel of models, one
can ask: given a new model's response pattern, what is the maximum
likelihood estimate of $\theta$, and what is its standard error?
This makes ``model X is at ability level $\theta_X$ on the
construct $\mathcal{B}$ measures'' a quantity with a confidence
interval, rather than a point score. It also lets us spot
benchmark-internal pathologies: items with very low $a_i$ (every
model gets them right or every model gets them wrong) carry no
discriminative signal and contribute only noise to the aggregate;
items with anomalously high $a_i$ flag potential contamination
(models that have seen them score sharply better than their
$\theta$ would otherwise predict).

\begin{intuition}
The 2PL form is the right baseline mental model for a single MCQ
benchmark. Imagine sweeping a model's $\theta$ across the real line
while holding $(a, b)$ fixed: at $\theta \ll b$, the model is below
the item's difficulty and answers near random; at $\theta \gg b$, the
model is above and answers near $1$; the transition width is set by
$a$. Frontier models on saturated MMLU items have $\theta_X$ well
above the items' $b_i$, so the response is near $1$ and provides no
new information --- this is the formal version of the saturation
intuition. Returning to canonical form: a saturated benchmark is
one where $\Pr[\text{correct} \mid \theta_X, b_i, a_i] \to 1$ for all
$i$, and $\hat{S}$ no longer responds to changes in $\theta_X$.
\end{intuition}

The IRT framing also clarifies what a successor benchmark like
MMLU-Pro \citep{wang2024-mmlu-pro} or HLE \citep{phan2025-hle} is
doing: shifting the item-difficulty distribution $\{b_i\}$ rightward
along the $\theta$ axis so frontier models are no longer in the
saturated regime, and selecting items with high $a_i$ that
discriminate among them. The construction recipes differ
(MMLU-Pro: $4 \to 10$-way MCQ plus reasoning-focused items; HLE:
expert-authored items with verifier discipline inherited from
GPQA and FrontierMath) but the abstract goal is the same: rebuild
the discriminative shoulder of the 2PL curve over the frontier-model
ability range. \deepencite{baker2001-irt-fundamentals §3, citation-pending}

\subsection{What the rest of the paper inherits}
\label{sec:measuring-handoff}

\Cref{sec:knowledge-benchmarks,sec:reasoning-benchmarks} take the
static-instance form $\mathcal{B} = (\rho, A, V)$ from
\eqref{eq:bench-static} and walk the lineage of difficulty
calibrations and verifier audits along the knowledge-vs-reasoning
boundary; \cref{sec:contamination} adds the contamination layer that
biases $\hat{S}$ in \eqref{eq:score-rv} when items leak;
\cref{sec:agentic-benchmarks} extends the static instance to the
agentic case. \Cref{sec:adversarial-robustness} switches the axis-3
coordinate from \emph{static} to \emph{adversarial} and re-instantiates
a benchmark family $(\mathcal{B}, \mathcal{A}, V, \pi)$ where
$\mathcal{A}$ is an attacker policy.
\Cref{sec:dangerous-capability} keeps axis-3 static but switches the
construct from refusal-rate to underlying capability for harm.

\Cref{sec:sycophancy,sec:scheming,sec:alignment-faking} take the
\emph{alignment} side of axis 1: their construct is not capability
but propensity to behave aligned, and the central methodological
problem is that their instruments measure capability under
elicitation, not propensity in deployment. Whether they admit a
deployment audit is the load-bearing open question for the
alignment-threat layer; the gap is the third coordinate of axis 3
still being under-instrumented in 2026.

\Cref{sec:watermarking} and \cref{sec:scalable-oversight} reach
beyond the in-distribution measurement frame entirely: watermarking
is a signature on generations rather than a score on a fixed set;
scalable oversight asks whether a bounded-resource judge can align a
model whose capability already exceeds the judge's. Both are
downstream of the vocabulary fixed here without being expressible
solely within it. We turn first to the knowledge layer
(\cref{sec:knowledge-benchmarks}) and the contamination boundary
(\cref{sec:contamination}) where the vocabulary is most directly
exercised.
